\documentclass[11pt]{article}
\usepackage[margin=1.05in]{geometry}
\usepackage{booktabs}
\usepackage{amsmath,amssymb}
\usepackage{microtype}
\usepackage{caption}
\usepackage{natbib}
\usepackage[colorlinks=true,linkcolor=black,citecolor=black,urlcolor=blue]{hyperref}
\captionsetup{font=small,labelfont=bf}
\bibliographystyle{plainnat}

\title{\vspace{-2.2em}What Makes a Good Fantasy Football Manager?}
\author{Evidence from one league over ten years and 252 leagues over one}
\date{September 2026}

\begin{document}
\maketitle
\vspace{-1.2em}

\begin{abstract}
\noindent Fantasy football managers hold two incompatible beliefs at once: that
the game is mostly luck, and that some managers are reliably good. We reconcile
them by decomposing weekly team scores into persistent manager quality and
week-to-week noise, and deriving the probability that the better manager wins.
We began expecting to find a low and general ceiling on skill. That is not what
the data show. In one fourteen-team league observed over four seasons the
intraclass correlation is $0.039$ --- a manager one standard deviation above
average wins only $55.7\%$ of his weeks, and $0.110$ estimated season by season
--- but across 252 public leagues the median is $0.294$, where the same manager
wins $67.6\%$. Skill content is
therefore a property of the \emph{league}, not of the game. About half the
difference is attributable to disengagement rather than judgment: dropping the
two lowest-scoring managers from each league moves the median from $0.294$ to
$0.164$, and the typical worst team scores $80\%$ of its league average against
$93\%$ in the tightly-contested league we study in detail. Within that league,
what most distinguishes managers is early-round draft outcomes ($r=0.805$ with
persistent scoring) --- but draft performance does not persist across eras
($r=-0.157$), so what looks like drafting skill is largely a draft that worked.
The one behaviour that is both learnable and demonstrably rewarded is waiver
claim volume.
\end{abstract}

\section{Introduction}

Ask a fantasy football league who its best manager is and you will get a
confident answer. Ask the same people whether the game is mostly luck and most
will say yes. Both convictions are usually held by the same person, and neither
is usually accompanied by a number.

They are not, in fact, contradictory. If persistent differences in manager
quality are small relative to week-to-week scoring noise, a real and stable
advantage still produces a win rate near a coin flip in any single week, and a
fourteen-game season is far too short to reveal it. Whether that describes a
given league is an empirical question with a specific answer, and the answer
turns out to vary enormously from league to league.

We set out to estimate the ceiling on manager skill, expecting to find it low
and roughly universal. The first half of that was right for the league we know
best and wrong in general. Our contributions are:

\begin{enumerate}
\item A variance decomposition of weekly scores into persistent manager quality
and noise, and a closed-form map from it to the probability that the better
manager wins (Section~\ref{sec:methods}).
\item Detailed estimates for one fourteen-team league over four seasons, where
$\rho = 0.039$ and the attainable edge is very small
(Section~\ref{sec:oneleague}).
\item The same estimate in 252 public leagues, where the median is $0.294$ and
the studied league sits near the 15th percentile on a like-for-like basis ---
falsifying the hypothesis that the ceiling is general (Section~\ref{sec:many}).
\item Evidence that roughly half of the typical league's apparent skill gap is
disengagement rather than judgment (Section~\ref{sec:dead}).
\item A within-league analysis of which manager behaviours track quality, and
the finding that the strongest one does not persist across a decade
(Sections~\ref{sec:behaviour}--\ref{sec:persistence}).
\end{enumerate}

\subsection*{Related work}

\citet{beckersun2016} formulate draft and lineup management as a mixed-integer
program, with weekly scoring modelled bilinearly as the product of a player's
innate ability at a statistic and his opponent's ability to defend it. Two of
their results anchor ours. Their Table~3 reports mean absolute weekly forecast
errors of $6.77$ points at running back against a mean of $7.46$, and $4.61$ at
tight end against $4.72$: the typical miss is 87--98\% of the typical score.
And when their optimizer is granted perfect knowledge of the season's outcomes,
it wins a ten-team league $24.7\%$ of the time against a $10\%$ baseline.
Omniscience is worth a factor of $2.5$.

\citet{lutz2015} applies support vector regression and neural networks to six
seasons of quarterback data, reaching a best mean absolute error of $6.22$
points and a mean relative error of $0.42$, without materially beating the
linear baselines he compares against.

Both papers ask how well weekly scores can be \emph{predicted}. We ask what the
irreducible residual implies for the attainable edge, and we work backwards from
realised outcomes rather than forwards from projections. We are not aware of a
published variance budget for head-to-head fantasy football, nor of an estimate
of how much that budget varies across leagues.

\section{Data}\label{sec:data}

\paragraph{The focal league.} A fourteen-team, half-PPR, head-to-head league
running since 2016, with a starting lineup of QB / RB / RB / WR / WR / TE /
FLEX / WRRB\_FLEX / K / DEF and four bench slots. It has used three platforms:
MyFantasyLeague 2016--19, Yahoo 2020--21, Sleeper since 2022. Primary analysis
uses the Sleeper era, from which we obtain every draft pick, every weekly matchup
with the full starter list and per-player points, every roster, and every
transaction \emph{including failed waiver claims with their bid amounts}. This
gives 770 manager-weeks, 798 draft picks, and 1{,}089 reconstructable waiver
auctions. We restrict to weeks 1--14, the fantasy regular season.

The earlier eras enter only in Section~\ref{sec:persistence}. Of those, only the
2016--17 drafts are recoverable: 2018 returns picks with an empty player field,
2019 returns no draft results, and the 2020--21 Yahoo export contains standings
and matchups but no player-level data. The MyFantasyLeague seasons are also
structurally different --- twenty franchises drafting as two \emph{independent}
ten-team drafts from a shared player pool --- and stored no scoring rules, so
player values there are computed in half-PPR and may not match what that league
actually scored.

\paragraph{The comparison sample.} To ask whether the focal league is typical we
snowball-sampled public Sleeper leagues, alternating between collecting the
leagues a known user belongs to and the members of a newly seen league, starting
from the focal league's members. We retained completed 2025 seasons with at
least eight managers and ten scored weeks, yielding \textbf{252 leagues}. Only
aggregate quantities are kept: league size, scoring type, roster positions, and
the matrix of weekly scores with managers reduced to integers. No usernames,
team names, or league names from third-party leagues are stored or reported.

This is a convenience sample reached through one league's social graph, and is
not a random draw from all Sleeper leagues. Section~\ref{sec:limits} returns to
what that does and does not threaten.

\section{Methods}\label{sec:methods}

\subsection{A season as an accounting identity}

Before any model, note that a season decomposes exactly. Let $O_m$ be the points
manager $m$'s \emph{optimal} lineup would have scored and $S_m$ what he actually
started. Every started player was either drafted by him or acquired later, so
\begin{equation}
O_m \;=\; \underbrace{D_m}_{\text{drafted starters}}
\;+\; \underbrace{A_m}_{\text{acquired starters}}
\;+\; \underbrace{L_m}_{\text{left on the bench}} ,
\label{eq:identity}
\end{equation}
with no residual. This is arithmetic, not estimation, and it is where we begin
because it fixes the scale of what follows.

\subsection{Variance components}

Let $Y_{mt}$ be the points scored by manager $m$'s starting lineup in week $t$.
We posit
\begin{equation}
Y_{mt} = \mu + \alpha_m + \varepsilon_{mt}, \qquad
\alpha_m \sim (0, \tau^2), \quad \varepsilon_{mt} \sim (0, \sigma^2),
\label{eq:model}
\end{equation}
with $\alpha_m \perp \varepsilon_{mt}$, and define the intraclass correlation
$\rho = \tau^2/(\tau^2+\sigma^2)$.

We stress what $\alpha_m$ is not. It is everything that persists week to week
--- including a draft that broke well, and including a manager who has stopped
setting his lineup. It is an upper bound on skill, not a measurement of it. Much
of this paper is about the gap between the two.

\subsection{From variance to win probability}

A matchup is decided by the margin
$D_{mm't} = (\alpha_m - \alpha_{m'}) + (\varepsilon_{mt} - \varepsilon_{m't})$.
The talent gap enters once; the noise enters \emph{twice}, with variance
$2\sigma^2$. Under approximate normality,
\begin{equation}
\Pr(m \text{ beats } m' \mid \delta) = \Phi\!\left(\frac{\delta}{\sigma\sqrt2}\right),
\qquad \delta = \alpha_m - \alpha_{m'} .
\end{equation}
Writing $\delta = c\tau$ for a manager $c$ standard deviations above average and
defining the \textbf{skill-to-noise ratio}
\begin{equation}
\kappa \;\equiv\; \frac{\tau}{\sigma\sqrt2} \;=\; \sqrt{\frac{\rho}{2(1-\rho)}},
\label{eq:kappa}
\end{equation}
this collapses to
\begin{equation}
\boxed{\;\Pr(\text{win a given week}) = \Phi(c\,\kappa)\;}
\label{eq:pk}
\end{equation}
The second equality in~\eqref{eq:kappa} follows from
$\tau^2 = \rho(\tau^2+\sigma^2)$ and $\sigma^2 = (1-\rho)(\tau^2+\sigma^2)$, and
is the reason everything downstream --- win rates, season records, the
observation horizon --- follows from $\rho$ alone. With $n$ independent
matchups, $W \sim \mathrm{Bin}(n, \Phi(c\kappa))$.

\subsection{Estimation}

Panels are unbalanced, so \eqref{eq:model} is fitted by unbalanced one-way ANOVA
\citep{searle1992}:
$\widehat\sigma^2 = \mathrm{MS}_W$ and
$\widehat\tau^2 = \max\{(\mathrm{MS}_B - \mathrm{MS}_W)/n_0,\,0\}$ with
$n_0 = (N - \sum_i n_i^2/N)/(k-1)$. The correction by $n_0$ is not cosmetic: in
the focal league the naive between-group variance of manager means is $26.2$, of
which $\widehat\sigma^2/\bar n = 9.5$ is sampling noise in the means themselves.
Omitting it overstates the persistent component by a third.

Intervals come from a nonparametric bootstrap resampling \emph{managers}, since
observations cluster within manager; $B=2000$.

\section{Results}

\subsection{One league, in detail}\label{sec:oneleague}

Applying~\eqref{eq:identity} to the focal league gives a first orientation:

\begin{center}\small
\begin{tabular}{lrr}
\toprule
& mean per season & between-manager SD \\
\midrule
$D$ \ points from drafted starters  & $1084$ & $207$ \\
$A$ \ points from acquired starters & $418$  & $161$ \\
$L$ \ points left on the bench      & $121$  & $41$ \\
\midrule
$O$ \ optimal total                 & $1623$ & \\
\bottomrule
\end{tabular}
\end{center}

Two-thirds of a manager's ceiling comes from his draft, a quarter from the wire,
and $7\%$ is forfeited to lineup mistakes.

Table~\ref{tab:budget} gives the variance decomposition.

\begin{table}[h]
\centering\small
\begin{tabular}{lrlrl}
\toprule
& \multicolumn{2}{c}{actual lineup} & \multicolumn{2}{c}{optimal lineup} \\
\cmidrule(lr){2-3}\cmidrule(lr){4-5}
& est. & 95\% CI & est. & 95\% CI \\
\midrule
$\sigma^2$ (within manager) & $499.5$ & $[450,\,550]$ & $520.0$ & $[480,\,561]$ \\
$\tau^2$ (between managers) & $20.3$ & $[1.2,\,36.9]$ & $20.6$ & $[0.0,\,41.7]$ \\
$\rho$ (ICC)                & $0.039$ & $[0.003,\,0.065]$ & $0.038$ & $[0.000,\,0.071]$ \\
$\Phi(\kappa)$              & $0.557$ & $[0.514,\,0.574]$ & $0.556$ & $[0.500,\,0.578]$ \\
\bottomrule
\end{tabular}
\caption{Focal league, $k=14$ managers, 770 manager-weeks, 2022--2025. The
optimal-lineup columns replace each manager's actual starters with the best
legal lineup available from his roster that week.}
\label{tab:budget}
\end{table}

In natural units $\widehat\sigma = 22.4$ points against $\widehat\tau = 4.5$:
\textbf{about 96\% of the variation in a weekly score is within-manager noise.}
The matchup margin carries a noise SD of $31.6$ points against a talent gap SD
of $6.4$.

The optimal-lineup columns deserve emphasis. Substituting every manager's
\emph{best possible} lineup for his actual one moves $\rho$ from $0.039$ to
$0.038$. Managers do differ in lineup efficiency --- the observed range is
$91.1\%$ to $94.1\%$ of optimal --- but relative to $\sigma$ the differences do
not register. Perfect hindsight start/sit decisions for everyone would not
change who wins this league.

Applying~\eqref{eq:pk}, a manager one standard deviation above average wins
$55.7\%$ of weeks, expects $7.79$ wins in fourteen, and finishes above $.500$
with probability $0.566$. By Spearman--Brown, a single season's record has
reliability $0.362$; reaching $0.80$ requires about 99 weeks, or seven seasons.

\subsection{The ceiling is not general}\label{sec:many}

We expected these figures to be roughly typical. They are not.

\begin{table}[h]
\centering\small
\begin{tabular}{lrrr}
\toprule
& median & IQR & 10th--90th pct \\
\midrule
$\rho$ (ICC)                 & $0.294$ & $[0.171,\,0.381]$ & $[0.080,\,0.456]$ \\
$\Phi(\kappa)$, P(+1 SD wins) & $0.676$ & $[0.626,\,0.710]$ & $[0.582,\,0.741]$ \\
\bottomrule
\end{tabular}
\caption{252 public Sleeper leagues, 2025 season. Mean $\sigma = 27.5$ points
per week; mean $\tau = 17.6$.}
\label{tab:many}
\end{table}

The median league has an ICC of $0.294$ and a competent manager there wins
$67.6\%$ of his weeks rather than $55.7\%$. Only $14\%$ of leagues fall below
$\rho = 0.10$.

The comparison requires care, because the two estimates are not computed over the
same span. Every league in the sample is observed for a single season, whereas
Table~\ref{tab:budget} pools four. Pooling treats a manager as one entity across
four entirely different rosters, which attenuates $\tau^2$; estimating the focal
league season by season instead gives $0.141$, $0.029$, $0.302$ and $0.078$ for
2022--25, with a median of $0.110$. The like-for-like figure is therefore
$0.110$, and \textbf{the focal league sits near the 15th percentile of the
sample} --- at the 10th if 2025 alone is used. Unusual, but less extreme than the
pooled estimate suggests.

That gap is itself informative. Within a single season $11\%$ of weekly variance
is between teams; across four seasons only $3.9\%$ is attributable to the
\emph{person}. Most of what separates teams within a year is that year's roster,
not a durable property of the manager --- the same conclusion
Section~\ref{sec:persistence} reaches from draft records.

Whatever the ceiling on skill is, it is not a constant of the game: the sample
spans $\rho = 0.080$ to $0.456$ between its 10th and 90th percentiles.

\subsection{Roughly half the difference is disengagement}\label{sec:dead}

Public leagues contain abandoned teams, and an abandoned team scores poorly every
week --- which is exactly what $\alpha_m$ is built to capture, without being
skill in any useful sense. To test this we recomputed each league's ICC after
removing its lowest-scoring managers.

\begin{center}\small
\begin{tabular}{lrr}
\toprule
sample & $n$ leagues & median $\rho$ \\
\midrule
all managers                & 252 & $0.294$ \\
drop lowest-scoring manager & 252 & $0.217$ \\
drop two lowest             & 252 & $0.164$ \\
\bottomrule
\end{tabular}
\end{center}

Removing two managers per league nearly halves the median ICC. The supporting
descriptive is stark: \textbf{the median worst team scores $80\%$ of its league
average, and the 10th percentile scores $62\%$.} In the focal league the worst
team scores $93\%$ --- every manager sits within $7\%$ of the mean.

Two conclusions follow. First, a substantial share of what a naive variance
decomposition calls ``manager skill'' in a public league is somebody who stopped
playing. Second, this does not explain the gap away. Applying the identical treatment to the
focal league moves it only from $0.039$ to $0.035$, because it has no disengaged
team to remove: the sample median under drop-two is $0.164$ against the focal
league's $0.035$, a ratio of $4.7$. That league is genuinely, unusually compressed, and
the compression is not an artifact of who is dropped.

\subsection{League design changes the skill content}

Scoring format is associated with the ICC in the direction one would predict
from the volatility of the underlying statistics: receptions are far more stable
week to week than touchdowns, so reception-heavy scoring should shrink $\sigma$
without much affecting $\tau$.

\begin{center}\small
\begin{tabular}{lrr}
\toprule
& $n$ leagues & median $\rho$ \\
\midrule
half-PPR & 48  & $0.224$ \\
full PPR & 201 & $0.299$ \\
\midrule
10 teams    & 35  & $0.197$ \\
11--12 teams & 193 & $0.310$ \\
13+ teams   & 20  & $0.342$ \\
\bottomrule
\end{tabular}
\end{center}

Three standard-scoring leagues are omitted from the upper panel as too few to
summarise. Correlation between league size and ICC is $+0.138$. We report these as
associations only: the sample is observational, formats are confounded with the
kind of league that adopts them, and the size relation runs opposite to the
focal league, which is the largest size class and the least noisy. The
practical reading is nonetheless that a commissioner has some control over how
much skill his league rewards.

\subsection{What separates managers within a league}\label{sec:behaviour}

Within the focal league we estimated each manager's $\alpha_m$ by empirical-Bayes
shrinkage and correlated it against six measured behaviours.

\begin{center}\small
\begin{tabular}{lr}
\toprule
behaviour & $r$ with $\widehat\alpha_m$ \\
\midrule
early-round draft value        & $+0.805$ \\
early-round bust avoidance     & $+0.625$ \\
lineup efficiency              & $+0.448$ \\
bid discipline                 & $+0.429$ \\
waiver efficiency (pts/\$)     & $+0.429$ \\
waiver claim volume            & $+0.413$ \\
\bottomrule
\end{tabular}
\end{center}

With $n=14$ managers and six tests this is exploratory. It also carries a
mechanical component we do not want to oversell: $\alpha_m$ is the persistent
part of weekly scoring, and the persistent part of a roster is mostly what was
drafted, so the leading correlation is close to arithmetic.

The waiver results are on much firmer ground, resting on 1{,}089 auctions rather
than sixteen picks per manager. Across the fourteen managers, claims attempted
per season correlates with waiver production at $r=+0.888$; price paid per
successful claim at $r=-0.799$; and claim \emph{win rate} at $r=-0.369$. The
negative sign on win rate is informative: a high win rate means bidding only when
certain, and certainty is expensive. Restricting to weeks 1--6 so that
acquisition date does not confound the comparison, players won for \$0 returned
$75.7$ points and players won for \$35 or more returned $78.4$ --- a $3.6\%$
difference for an unbounded price, with points per dollar falling from $29.9$ in
the \$1--4 band to $1.7$ above \$35. Since FAAB expires worthless and bench slots
recycle weekly, the marginal cost of an additional small claim is near zero.

We note one correction that changed this section materially. Our first pass
credited a claimed player with all points he scored while rostered.
Recomputing against weeks he was actually \emph{started} reduced every manager's
waiver production by a median of $65\%$, and non-uniformly ($49\%$ to $78\%$
across managers), which reordered the table.

\subsection{Draft performance does not persist}\label{sec:persistence}

If drafting is a stable trait, it should correlate across time. For the twelve
managers present in both eras, the correlation between standardised draft score
in 2016--17 and in 2022--25 is $r = -0.157$, 95\% CI $[-0.67,\,+0.46]$ by
Fisher's $z$. Eight of twelve changed sign; the manager with the best four-year
draft record in the current era had the worst in the earlier one.

This estimate is attenuated by measurement error in both halves. With within-era
ICC $0.215$ for season draft scores, Spearman--Brown gives reliabilities of
$0.523$ and $0.354$ for four- and two-season means, so the disattenuated
correlation is $-0.157/\sqrt{0.523 \times 0.354} = -0.365$. Correcting for
attenuation moves the estimate \emph{further} from zero in the negative
direction; it cannot rescue a claim of persistence. Supporting
$r_{\text{true}} = +0.50$ would have required observing $r_{\text{obs}} = +0.215$.

Pooling all six gradeable seasons, the share of variance in season draft scores
attributable to the manager falls from $0.215$ to $0.027$.

We hold this loosely. The eras differ in league size, roster structure and
possibly scoring, and four years separate them. It is evidence of absence only
in the weak sense: we could not detect persistence with the data available.

\subsection{What is outside a manager's control}

Two factors commonly blamed or credited can be measured directly. Draft slot is
randomly assigned; across 56 manager-seasons its correlation with season wins is
$+0.049$. Schedule is not: regressing wins on points scored gives $R^2 = 0.484$,
and adding points \emph{allowed} raises it to $0.719$. \textbf{Roughly a quarter
of the variance in season records is explained purely by whom a manager happened
to be scheduled against.}

\subsection{Validation, and where it fails}\label{sec:validation}

The model is fitted to weekly scores and makes a testable prediction about season
records. If win rates are $\mathrm{Bin}(14,p_m)/14$ with $p_m = \Phi(c_m\kappa)$,
the talent-induced SD of the win rate is approximately $2\phi(0)\kappa$.

\begin{center}\small
\begin{tabular}{lr}
\toprule
observed SD of season win rate & $0.1517$ \\
SD implied by a pure coin flip, $n=14$ & $0.1336$ \\
residual attributable to talent & $0.0718$ \\
predicted from the weekly model & $0.1137$ \\
\bottomrule
\end{tabular}
\end{center}

\textbf{These do not agree.} The weekly model predicts a talent spread in season
records $1.6$ times larger than the records display. The discrepancy is not
fatal --- $\widehat\tau^2$ has interval $[1.2,\,36.9]$, and $0.0718$ corresponds
to a $\tau$ near its lower end --- but the point estimate in
Table~\ref{tab:budget} should be read as an upper bound. Candidate explanations
are non-normality in the margin, which heavier tails would produce, and
schedule-induced correlation in opponent quality.

\section{Discussion}

\subsection{Repetition, not insight}

Table~\ref{tab:reps} organises the decisions a manager makes by how often he
makes them.

\begin{table}[h]
\centering\small
\begin{tabular}{lrr}
\toprule
decision & per season & between-manager SD of season total \\
\midrule
trades                 & $\sim1$  & $\sim30$ \\
early draft picks      & $4$      & $207$ \\
weekly lineups         & $14$     & $41$ \\
waiver claims (attempted) & $32$ & $161$ \\
\bottomrule
\end{tabular}
\caption{Focal league. Waiver claims are attempts including failed bids, median
per manager-season (range 7--92); completed claims median 15. Trades are rare and
shrinking: nine in 2022, two in 2025.}
\label{tab:reps}
\end{table}

A per-decision edge accumulates linearly in the number of decisions, while the
noise in \emph{measuring} that edge grows only as $\sqrt n$. Repetition thus
improves both how much advantage a manager can bank and our ability to detect
that he has it. This organises the paper's otherwise scattered findings. Waivers
produce the cleanest signal in the data because a manager attempts about thirty a
season and the sample contains 1{,}089 auctions. Draft skill does not survive an era change partly
because four meaningful picks a year over four years is sixteen noisy
observations. Lineup selection never registers because its per-decision effect
--- about $2.8$ points a week --- is an order of magnitude below the $31.6$
points of matchup noise it competes with.

\subsection{Two different questions}

``What makes a good fantasy manager'' turns out to have two answers that are
routinely conflated.

\emph{What makes a manager score more this season?} Predominantly a draft that
worked ($r = 0.805$). That is largely outside his control and does not carry from
one era to the next.

\emph{What can a manager actually improve?} Waiver behaviour: more claims, at
lower prices. Its correlation with quality is weaker ($r = 0.413$), but its
feedback loop is fast enough to learn from within a single season.

Confusing the first for the second is, we suspect, why so many managers believe
they are good at drafting.

\section{Limitations}\label{sec:limits}

\paragraph{The comparison sample is not random.} It was reached by snowball
sampling from one league's social graph, so leagues socially proximate to the
focal league are over-represented, and $201$ of $252$ are full PPR. This
threatens the point estimate of the median ICC. It does not threaten the paper's
central claim, which is about \emph{dispersion}: whatever the population median,
we observe ICCs from $0.079$ to $0.456$ within a single sample, and that range
is what falsifies a universal ceiling.

\paragraph{Disengagement is inferred, not observed.} We identify likely
abandoned teams by low scoring, because the crawl retained only weekly scores.
Separating ``quit'' from ``bad but trying'' requires transaction counts per
league, which we did not collect. The claim that roughly half the typical
league's skill gap is disengagement should be read as consistent with the
evidence rather than established by it.

\paragraph{$\alpha_m$ is not skill.} It is everything persistent. We have no
design that decomposes an individual manager's $\alpha_m$ into judgment and
fortune; the cross-era comparison in Section~\ref{sec:persistence} is the closest
available and it is underpowered.

\paragraph{Normality.} Equation~\eqref{eq:pk} assumes an approximately normal
margin. Weekly scores are right-skewed sums of skewed player scores, and the
validation failure in Section~\ref{sec:validation} is what fat tails would
produce.

\paragraph{Free parameters and leakage.} The $-50$ bust threshold and the
$\pm12$ window in the slot-expectation curve were chosen without formal
justification; the ordering of managers is stable across thresholds from $-100$
to $0$ but the magnitude is not. The slot curve is fitted on the same picks it
scores, inducing mild leakage. Replacement level is computed from season totals
although lineups are set weekly, which inflates all points-above-replacement
figures.

\paragraph{Single-season ICCs are noisy.} Each comparison league contributes one
season, so $\tau$ there is estimated from roughly fourteen weeks per manager. The
focal league's own per-season estimates range from $0.029$ to $0.302$, which is
the scale of instability to expect in any individual league's figure. Conclusions
are drawn from the distribution across leagues, not from any single one.

\paragraph{Estimation.} Method-of-moments $\widehat\tau^2$ is truncated at zero,
and the bootstrap inherits that truncation --- visible in
Table~\ref{tab:budget}, where the optimal-lineup interval is pinned at the
boundary. A hierarchical Bayesian model with $\tau$ on the log scale, partial
pooling across leagues, and league covariates entering as coefficients would
estimate this properly and is the obvious next step.

\section{Conclusion}

The question that titles this paper has a smaller answer than most managers
believe, and a more variable one than we expected. In a competitive, fully
engaged fourteen-team league, a manager one standard deviation above average wins
$55.7\%$ of his weeks and would need seven seasons to demonstrate he is not
lucky. In the median public league the same manager wins $67.6\%$ --- but about
half of that larger edge appears to come from playing against people who have
stopped trying, and the comparison narrows further once both sides are estimated
over the same span.

For a manager wanting to improve, the practical implication is narrow and
somewhat deflating. The decision that most determines how good he looks --- the
first four rounds of his draft --- is the one he gets four attempts at per year
and cannot demonstrably learn on a human timescale. The decision he can actually
improve is the one that feels like administrative work: many small waiver claims
rather than few expensive ones. And a season's standings, in a league where
everybody is trying, are close to uninformative about which of the two he did
well.

\begin{thebibliography}{9}
\bibitem[Becker and Sun(2016)]{beckersun2016}
Becker, A. and X.~A. Sun (2016).
\newblock An analytical approach for fantasy football draft and lineup management.
\newblock \emph{Journal of Quantitative Analysis in Sports} 12(1), 17--30.

\bibitem[Lutz(2015)]{lutz2015}
Lutz, R. (2015).
\newblock Fantasy football prediction.
\newblock \emph{arXiv preprint} arXiv:1505.06918.

\bibitem[Searle et al.(1992)]{searle1992}
Searle, S.~R., G.~Casella, and C.~E. McCulloch (1992).
\newblock \emph{Variance Components}.
\newblock Wiley, New York.
\end{thebibliography}

\end{document}
